\documentclass[../main.tex]{subfiles}

\begin{document}

\chapter*{全书小结}

本书从大量重复试验中事件出现频率的稳定性这一经验事实出发，引进了一系列概念，建立了既有广泛实际应用又有深刻理论结果的一整套数学理论；最后这个理论又令人信服地解释了作为出发点的经验事实，到此为止，这个课程已告一段落。

读者可以看到，作为概率论研究对象的随机现象，在自然界和人类社会中是普遍存在的，这既说明了概率论理论的重要性，也决定了它的应用的广泛性。

读者也可以看到，概率论是数学的一个有特色的分支。一方面，由于它与其他数学分支（它们都研究决定性现象）研究对象的不同，因此它有着别开生面的研究课题，从而也有着自己独特的概念与方法。另一方面，它又是一个严谨的数学分支，它的概念有明确的定义，它的方法是严格的，它的结果是深刻的，这是许多著名数学家长期耕耘的结果，也得益于不断地从其他数学分支吸取有用的概念与方法。

把概率论最基本的概念——事件，概率，随机变量与数学期望分别看作是集合，规范化测度，可测函数与可测函数关于规范化测度的积分，这种观点自从概率论的公理化结构体系出现之后，已被普遍接受。这里强调的是概率论的测度论基础，对于概率论基本概念的明确定义是至关重要的。

读者应当认识到，上述观点的形成是一个历史过程，同时也不应当忘记这些概念的现实背景，这时研究几个经典的模型是很有帮助的。历史上，古典概型、几何概率、Bernoulli概型等几个模型孕育了早期的概率论，对后来的发展也有重大影响。即使在今天，这些既直观又具体的模型也还值得我们特别注意。

统计独立性是概率论中特有的概念，它的引进大大丰富了概率论的研究。概率论基础中最深入的结果大都是在独立性的假定下获得的，这主要是指几种形式的极限定理：大数定律、强大数定律与中心极限定理。

条件概率是另一重要概念，它使我们能充分利用有关的信息，在概率计算中十分有用。当进一步研究非独立的场合时，条件概率将起更大的作用。

随机变量概念的普遍使用相对来讲还是近代的事，但是这个概念是很自然的，又便于运算，用它来描述事件有不少好处。涉及随机变量的最重要问题中大多数可以通过分布函数来表述，概率论的这一部分可以独立于它的测度论基础而进行，这里所用的主要是分析方法。

一种分布就是一个数学模型，每种分布有自己的特征；分布之间有各种联系。书中把分布按重要性分成三类：第一类三大分布有专门的节加以介绍；第二类包括十几种重要分布，在正文中指明其背景、性质以及它与其他分布的关系；第三类则多数在例题或习题中出现。最后它们大多在附录四中汇总。

数学期望是概率论中最古老的概念，它的明确的直观含意和良好的数学性质使它在概率论中一直占据着重要的地位，各种重要的数字特征大都是某种数学期望，甚至概率也是一种数学期望。数学期望，以及与它相辅相成的方差，刻画了随机变量的概貌，是实际应用中最关心的两个量。

在多维场合，随机向量由多元联合分布完整描述，用均值向量和协方差矩阵刻画概貌，而边际分布、条件分布、独立性和相关系数等概念更是丰富了研究主题。

极限定理是概率论中最重要的理论结果，本书通过模块式的三个层次介绍了经典结果，相当详尽。主要假设是独立性和矩的存在，而分布函数、矩和特征函数这三者则是解决古典极限定理的主要工具。

总的说来，本书通过经典模型来提供有关背景，采用公理化结构以明确定义概念，强调独立性以突出学科的特点，利用分析方法来获得深刻的结果，最后也是最重要的一点是，试图通过大量的实例来介绍概率论的日益广泛的应用。

正如本书的书名所指出的，这些知识仅仅是概率论的基础，它们是新的研究的出发点。

\end{document}